\item \model{مرور جامع ادبیات پژوهش: تقطیر دانش و فشرده‌سازی در مدل‌های زبانی بزرگ}

\paragraph*{مقدمه و سیر تکامل تاریخی تقطیر در هوش مصنوعی}
با افزایش تصاعدی مقیاس مدل‌های زبانی بزرگ (\en{LLMs}) به صدها میلیارد پارامتر، چالش‌های ناظر بر حافظه پردازشی، توان محاسباتی و هزینه‌های زیست‌محیطی در مراحل آموزش و استنتاج (\en{Inference})، تقطیر دانش (\en{Knowledge Distillation - KD}) را از یک راهکار فرعی به یکی از ارکان بنیادی مهندسی هوش مصنوعی بدل ساخته است \cite{mansourian2025comprehensive, fang2026knowledge, song2026survey}. تقطیر دانش که نخستین بار توسط بوچیلوآ و همکاران \cite{bucilua2006model} برای مدل‌های بیزین و سپس توسط هینتون و همکاران \cite{hinton2015distilling} در یادگیری عمیق صورتبندی شد، فرآیندی است که طی آن یک مدل فشرده و سبک تحت عنوان دانش‌آموز (\en{Student}) برای بازتولید رفتار و بازنمایی‌های درونی یک مدل حجیم تحت عنوان آموزگار (\en{Teacher}) آموزش می‌بیند \cite{wen2025knowledge}.

در طول یک دهه اخیر، مفهوم تقطیر از تطبیق ایستا و نقطه‌ای لاجیت‌ها در مسائل طبقه‌بندی تصویر به حوزه‌های بسیار پیچیده‌ای نظیر استدلال زنجیره تفکر (\en{Chain-of-Thought - CoT})، تقطیر درون‌خط‌مشی (\en{On-Policy Distillation - OPD})، هم‌افزایی با تقطیر مجموعه داده (\en{Dataset Distillation - DD}) و کوانتیزاسیون مجانبی بدون اتلاف آماری (\en{Statistically-Lossless Quantization}) تکامل یافته است.

\paragraph*{مبانی تحلیلی تقطیر کلاسیک و مشتق‌گیری ریاضی رفتار مجانبی}
در صورتبندی پایه تقطیر دانش برای بردار لاجیت‌های آموزگار $z^T \in \mathbb{R}^{|\mathcal{V}|}$ و دانش‌آموز $z^S \in \mathbb{R}^{|\mathcal{V}|}$، توزیع‌های احتمال هموارشده از طریق تابع بیشینه هموار با دمای $\tau > 0$ به صورت زیر تعریف می‌شوند:
\begin{equation}
    p_i(z, \tau) = \frac{\exp(z_i / \tau)}{\sum_{j=1}^{|\mathcal{V}|} \exp(z_j / \tau)}
\end{equation}
تابع هزینه تلفیقی تقطیر بر پایه واگرایی کولبک-لایبلر (\en{Kullback-Leibler Divergence}) به فرم زیر است \cite{hinton2015distilling, mansourian2025comprehensive}:
\begin{equation}
    \mathcal{L}_{\text{KD}} = \alpha \mathcal{L}_{\text{CE}}(y, p(z^S, 1)) + (1 - \alpha) \tau^2 \mathcal{D}_{\text{KL}}\left(p(z^T, \tau) \,\|\, p(z^S, \tau)\right)
\end{equation}
که در آن:
\begin{equation}
    \mathcal{D}_{\text{KL}}\left(p(z^T, \tau) \,\|\, p(z^S, \tau)\right) = \sum_{i=1}^{|\mathcal{V}|} p_i(z^T, \tau) \log \left( \frac{p_i(z^T, \tau)}{p_i(z^S, \tau)} \right)
\end{equation}

جهت بررسی سازوکار انتشار گرادیان، مشتق جزئی واگرایی $\mathcal{D}_{\text{KL}}$ را نسبت به لاجیت دانش‌آموز $z_i^S$ محاسبه می‌کنیم. با بسط رابطه داریم:
\begin{equation}
    \frac{\partial \mathcal{D}_{\text{KL}}}{\partial z_i^S} = - \sum_{j=1}^{|\mathcal{V}|} p_j(z^T, \tau) \frac{\partial \log p_j(z^S, \tau)}{\partial z_i^S} = - \sum_{j=1}^{|\mathcal{V}|} \frac{p_j(z^T, \tau)}{p_j(z^S, \tau)} \frac{\partial p_j(z^S, \tau)}{\partial z_i^S}
\end{equation}
با توجه به مشتق تابع سافت‌مکس که برابر است با:
\begin{equation}
    \frac{\partial p_j(z^S, \tau)}{\partial z_i^S} = \frac{1}{\tau} p_i(z^S, \tau) (\delta_{ij} - p_j(z^S, \tau))
\end{equation}
با جایگذاری رابطه فوق در گرادیان واگرایی خواهیم داشت:
\begin{equation}
    \frac{\partial \mathcal{D}_{\text{KL}}}{\partial z_i^S} = - \sum_{j=1}^{|\mathcal{V}|} \frac{p_j(z^T, \tau)}{p_j(z^S, \tau)} \left[ \frac{1}{\tau} p_i(z^S, \tau) (\delta_{ij} - p_j(z^S, \tau)) \right] = \frac{1}{\tau} \left( p_i(z^S, \tau) - p_i(z^T, \tau) \right)
\end{equation}

اکنون رفتار این تابع را در حد دمای بسیار بالا ($\tau \to \infty$) با استفاده از بسط تیلور بررسی می‌کنیم. بسط نمایی حول صفر عبارت است از:
\begin{equation}
    \exp\left(\frac{z_i}{\tau}\right) = 1 + \frac{z_i}{\tau} + \frac{z_i^2}{2\tau^2} + \mathcal{O}\left(\frac{1}{\tau^3}\right)
\end{equation}
با فرض میانگین-صفر بودن لاجیت‌ها بر روی کل واژگان ($\sum_j z_j = 0$):
\begin{equation}
    \sum_{j=1}^{|\mathcal{V}|} \exp\left(\frac{z_j}{\tau}\right) \approx |\mathcal{V}| + \frac{1}{\tau}\sum_{j=1}^{|\mathcal{V}|} z_j = |\mathcal{V}|
\end{equation}
در نتیجه، احتمال نرمال‌شده به صورت زیر تقریب زده می‌شود:
\begin{equation}
    p_i(z, \tau) \approx \frac{1}{|\mathcal{V}|} + \frac{z_i}{|\mathcal{V}|\tau}
\end{equation}
با قرار دادن این تقریب در رابطه گرادیان، اثر ضریب $\tau^2$ در گرادیان کل مشخص می‌گردد:
\begin{equation}
    \tau^2 \frac{\partial \mathcal{D}_{\text{KL}}}{\partial z_i^S} \approx \tau^2 \cdot \frac{1}{\tau} \left( \left[\frac{1}{|\mathcal{V}|} + \frac{z_i^S}{|\mathcal{V}|\tau}\right] - \left[\frac{1}{|\mathcal{V}|} + \frac{z_i^T}{|\mathcal{V}|\tau}\right] \right) = \frac{1}{|\mathcal{V}|} (z_i^S - z_i^T)
\end{equation}
این نتیجه تحلیلی اثبات می‌کند که در حد دمای بالا، تقطیر بر پایه واگرایی \en{KL} به طور کامل معادل با **کمینه‌سازی خطای میانگین مربعات (\en{MSE})** بین لاجیت‌های میانگین-صفر دو مدل است:
\begin{equation}
    \lim_{\tau \to \infty} \tau^2 \mathcal{D}_{\text{KL}}\left(p(z^T, \tau) \,\|\, p(z^S, \tau)\right) = \frac{1}{2|\mathcal{V}|} \|z^S - z^T\|_2^2
\end{equation}

\paragraph*{رده‌بندی پنج‌گانه روش‌های تقطیر دانش}
طبقه‌بندی جامع ارائه‌شده توسط منصوریان و همکاران \cite{mansourian2025comprehensive} نشان می‌دهد که تقطیر مدرن بر پنج محور اساسی استوار است که در نمودار زیر ساختاربندی شده‌اند:

\begin{center}
\begin{tikzpicture}[
    box/.style={rectangle, draw=blue!70!black, fill=blue!5, rounded corners=3pt, text centered, font=\small\bfseries, minimum height=0.8cm, inner sep=4pt},
    leaf/.style={rectangle, draw=gray!60, fill=white, rounded corners=2pt, text centered, font=\footnotesize, minimum height=0.7cm, inner sep=3pt},
    arrow/.style={-{Stealth[scale=0.8]}, draw=blue!60!black, thick}
]
    \node[box, fill=blue!20] (root) at (0, 0) {ابعاد بنیادین تقطیر دانش (\en{KD})};
    
    \node[box] (src) at (-5.5, -1.8) {منابع دانش (\en{Sources})};
    \node[box] (sch) at (-2.7, -1.8) {طرح‌ها (\en{Schemes})};
    \node[box] (alg) at (0, -1.8) {الگوریتم‌ها (\en{Algorithms})};
    \node[box] (mod) at (2.8, -1.8) {روش‌های ورودی (\en{Modalities})};
    \node[box] (app) at (5.7, -1.8) {کاربردها (\en{Applications})};
    
    \draw[arrow] (root) -- (src);
    \draw[arrow] (root) -- (sch);
    \draw[arrow] (root) -- (alg);
    \draw[arrow] (root) -- (mod);
    \draw[arrow] (root) -- (app);
    
    \node[leaf, align=center] (src_c) at (-5.5, -3.2) {• لاجیت‌ها (\en{Logits})\\• ویژگی‌ها (\en{Features})\\• شباهت‌ها (\en{Similarities})};
    \node[leaf, align=center] (sch_c) at (-2.7, -3.2) {• آفلاین (\en{Offline})\\• آنلاین (\en{Online})\\• خود-تقطیری (\en{Self})};
    \node[leaf, align=center] (alg_c) at (0, -3.2) {• مبتنی بر توجه\\• تخاصمی (\en{GAN})\\• تقطیر تطبیقی\\• متقابل (\en{Contrastive})};
    \node[leaf, align=center] (mod_c) at (2.8, -3.2) {• متنی (\en{NLP})\\• بینایی و تصویر\\• چندوجهی و ابرنقاط\\• صوتی (\en{Speech})};
    \node[leaf, align=center] (app_c) at (5.7, -3.2) {• مدل‌های زبانی\\• مدل‌های پایه (\en{CLIP})\\• بینایی خودرویی\\• بیوانفورماتیک};
    
    \draw[draw=gray!50, dashed] (src) -- (src_c);
    \draw[draw=gray!50, dashed] (sch) -- (sch_c);
    \draw[draw=gray!50, dashed] (alg) -- (alg_c);
    \draw[draw=gray!50, dashed] (mod) -- (mod_c);
    \draw[draw=gray!50, dashed] (app) -- (app_c);
\end{tikzpicture}
\end{center}

\begin{itemize}
    \item \textbf{تقطیر مبتنی بر لاجیت پیشرفته:} تفکیک توزیع کلاس هدف و غیرهدف در روش \en{DKD} \cite{zhao2022decoupled}، استانداردسازی پیش‌پردازش لاجیت در \en{NormKD} \cite{chi2023normkd} و نگاشت کروی برای خنثی‌سازی خطای مقیاس در \en{SphericalKD} \cite{guo2020reducing}.
    \item \textbf{تقطیر مبتنی بر ویژگی (\en{Feature-based}):} نگاشت حالت‌های پنهان میانی با تبدیل‌های خطی در \en{FitNet} \cite{romero2014fitnets}، تطبیق توجه در \en{AT} \cite{zagoruyko2016paying}، تقطیر نرمال‌سازی کانالی در \en{CWD} \cite{shu2021channel} و کاربرد مدل‌های دیفیوژن در پالایش ویژگی‌ها در روش \en{DiffKD} \cite{huang2023knowledge}.
    \item \textbf{تقطیر مبتنی بر شباهت (\en{Similarity-based}):} انتقال فاصله اقلیدسی و زاویه‌ای میان نمونه‌ها در \en{RKD} \cite{park2019relational}، همبستگی کانالی در \en{ICKD} \cite{liu2021exploring} و ساختار رابطه‌ای بین‌کلاسی در \en{DistKD} \cite{huang2022knowledge} و \en{AICSD} \cite{mansourian2025aicsd}.
\end{itemize}

\paragraph*{تقطیر درون‌خط‌مشی (\en{On-Policy Distillation - OPD}) و حذف خطای انباشته مواجهه}
در استنتاج خودبازگشتی مدل‌های زبانی، ناتوانی روش‌های برون‌خط‌مشی (\en{Off-Policy}) ناشی از پدیده **تورش مواجهه (\en{Exposure Bias})** است \cite{song2026survey, wen2025knowledge}. در روش‌های مرسوم، مدل روی پیشوندهای کاملاً صحیح تولیدشده توسط انسان یا آموزگار ($y_{<t} \sim p_{\text{data}}$) آموزش می‌بیند، در حالی که در استنتاج، مدل از خروجی‌های خود استفاده می‌کند ($y_{<t} \sim p_\theta$).

\textbf{تحلیل قضیه دگر (\en{DAgger Compounding Error Bound}):} طبق نتایج راس و بگنل \cite{ross2011reduction}، چنانچه احتمال خطای پیش‌بینی در هر گام $\epsilon$ باشد، خطای کل در یک افق زمانی $T$ توکنی در آموزش ایستا (برون‌خط‌مشی) به صورت درجه دوم رشد می‌کند:
\begin{equation}
    \mathbb{E}[\text{Error}]_{\text{Off-Policy}} = \mathcal{O}(\epsilon T^2)
\end{equation}
زیرا بروز یک خطای منفرد در ابتدای دنباله، مدل را به مناطقی از فضای توزیع می‌برد که هرگز در فاز آموزش مشاهده نکرده است. در مقابل، تقطیر درون‌خط‌مشی (\en{OPD}) از طریق وادار ساختن دانش‌آموز به نمونه‌برداری از توزیع خود ($y \sim p_\theta$) و دریافت فیدبک نظارتی روی همین مسیرها، نرخ تجمع خطا را به مرتبه خطی بازمی‌گرداند:
\begin{equation}
    \mathbb{E}[\text{Error}]_{\text{On-Policy}} = \mathcal{O}(\epsilon T)
\end{equation}

\paragraph*{چارچوب یکپارچه واگرایی‌های $f$ (\en{Unified $f$-Divergence Framework})}
سونگ و ژنگ \cite{song2026survey} نشان دادند که رویکردهای نوین تقطیر درون‌خط‌مشی، همگی نمونه‌هایی خاص از یک چارچوب بهینه‌سازی واگرایی $f$ روی توزیع نمونه‌برداری خط‌مشی دانش‌آموز هستند:
\begin{equation}
    \mathcal{L}_{\text{OPD}}(\theta) = \mathbb{E}_{x \sim \mathcal{D}, \, y \sim p_\theta(\cdot \mid x)} \left[ \sum_{t=1}^{|y|} \mathcal{D}_f\left(p_T(\cdot \mid x, y_{<t}) \,\|\, p_\theta(\cdot \mid x, y_{<t})\right) \right]
\end{equation}
که در آن عملگر واگرایی برای تابع مولد محدب $f$ با شرط $f(1)=0$ به صورت زیر تعریف می‌شود:
\begin{equation}
    \mathcal{D}_f(P \,\|\, Q) = \mathbb{E}_{y \sim Q} \left[ f\left(\frac{P(y)}{Q(y)}\right) \right]
\end{equation}
ویژگی‌های هندسی واگرایی‌ها در برازش توزیع‌های زبانی عبارتند از:
\begin{itemize}
    \item \textbf{واگرایی مستقیم (\en{Forward KL}):} با $f(u) = u \log u$ منجر به رفتار **پوشش مد (\en{Mode-Covering})** می‌شود. دانش‌آموز تمامی مدهای آموزگار را در بر می‌گیرد و در مسائل دارای پاسخ‌های متکثر (نظیر ترجمه یا خلاصه‌سازی) از افت تنوع جلوگیری می‌کند، هرچند ممکن است به توهم (\en{Hallucination}) در فواصل بین مدی بیانجامد.
    \item \textbf{واگرایی معکوس (\en{Reverse KL}):} با $f(u) = -\log u$ منجر به رفتار **جستجوی مد (\en{Mode-Seeking})** می‌شود. دانش‌آموز تنها بر یک قله با احتمال بالا متمرکز شده و از ورود به نواحی با چگالی ناچیز آموزگار شدیداً اجتناب می‌ورزد؛ این راهبرد برای مسائل استدلال ریاضیاتی، فرمول‌ها و کدنویسی که پاسخ واحد دارند ایده‌آل است.
    \item \textbf{واگرایی‌های آمیخته و چوله (\en{Skewed KL}):} روش \en{DistiLLM} \cite{ko2024distillm} با هدف پیشگیری از ناپایداری‌های عددی و تقسیم بر صفر هنگامی که $p_T(y) \approx 0$ یا $p_\theta(y) \approx 0$ است، توزیع چوله $\tilde{p} = \alpha p_T + (1 - \alpha) p_\theta$ را تعبیه کرده و واگرایی $\mathcal{D}_{\text{KL}}(p_T \,\|\, \tilde{p})$ را کمینه می‌سازد.
\end{itemize}

\paragraph*{اثبات هم‌ارزی تحلیلی \en{Reverse KL} با گرادیان خط‌مشی در \en{RL}}
در روش‌های تقطیر در سطح دنباله نظیر \model{MiniLLM} \cite{gu2024minillm}، هدف کمینه‌سازی واگرایی معکوس روی دنباله کامل خروجی است:
\begin{equation}
    \mathcal{L}_{\text{Seq-RKL}}(\theta) = \mathbb{E}_{y \sim p_\theta(\cdot \mid x)} \left[ \log p_\theta(y \mid x) - \log p_T(y \mid x) \right]
\end{equation}
گرادیان این تابع هدف نسبت به $\theta$ مستقیماً بر مبنای قضیه گرادیان خط‌مشی به دست می‌آید:
\begin{equation}
    \nabla_\theta \mathcal{L}_{\text{Seq-RKL}}(\theta) = \nabla_\theta \sum_{y} p_\theta(y \mid x) \left[ \log p_\theta(y \mid x) - \log p_T(y \mid x) \right]
\end{equation}
با بسط مشتق و استفاده از قاعده لایبنیتس:
\begin{equation}
    \nabla_\theta \mathcal{L} = \sum_{y} \left( \nabla_\theta p_\theta(y \mid x) \right) \left[ \log \frac{p_\theta(y \mid x)}{p_T(y \mid x)} \right] + \sum_{y} p_\theta(y \mid x) \nabla_\theta \left[ \log p_\theta(y \mid x) - \log p_T(y \mid x) \right]
\end{equation}
جمله دوم به دلیل نرمال بودن توزیع احتمال صفر می‌شود:
\begin{equation}
    \sum_{y} p_\theta(y \mid x) \nabla_\theta \log p_\theta(y \mid x) = \sum_y \nabla_\theta p_\theta(y \mid x) = \nabla_\theta \left(\sum_y p_\theta(y \mid x)\right) = \nabla_\theta(1) = 0
\end{equation}
با اعمال ترفند نسبت درست‌نمایی ($\nabla_\theta p_\theta = p_\theta \nabla_\theta \log p_\theta$) روی جمله اول، رابطه زیر به دست می‌آید:
\begin{equation}
    \nabla_\theta \mathcal{L}_{\text{Seq-RKL}}(\theta) = - \mathbb{E}_{y \sim p_\theta} \left[ \left( \log \frac{p_T(y \mid x)}{p_\theta(y \mid x)} \right) \nabla_\theta \log p_\theta(y \mid x) \right]
\end{equation}
اگر نسبت پاداش توکنی به صورت $r(x, y) = \log \frac{p_T(y \mid x)}{p_\theta(y \mid x)}$ تعریف گردد، این رابطه دقیقاً معادل با **فرمول کلاسیک \en{REINFORCE} در یادگیری تقویتی** خواهد بود \cite{gu2024minillm, song2026survey}.

\paragraph*{روش‌های تقطیر هدایت‌شده با یادگیری تقویتی و خود-تقطیری ممتاز (\en{OPSD})}
\begin{itemize}
    \item \textbf{برون‌یابی پاداش با \en{G-OPD}:} یانگ و همکاران \cite{yang2026learning} نشان دادند که می‌توان با بازتعریف فرآیند تقطیر به صورت بهینه‌سازی تحت قید \en{KL}:
    \begin{equation}
        \max_\theta \mathbb{E}_{y \sim p_\theta} \left[ \sum_{t=1}^{|y|} \alpha \log \frac{p_T(y_t \mid y_{<t})}{p_{\text{ref}}(y_t \mid y_{<t})} - \mathcal{D}_{\text{KL}}\left(p_\theta(\cdot \mid y_{<t}) \,\|\, p_{\text{ref}}(\cdot \mid y_{<t})\right) \right]
    \end{equation}
    به ازای $\alpha > 1$ پدیده برون‌یابی پاداش (\en{Reward Extrapolation}) را برانگیخت که طی آن دانش‌آموز استدلال‌هایی نوین و کوتاه‌تر خلق می‌کند که حتی ورای گستره آموزشی آموزگار قرار دارند.
    \item \textbf{خود-تقطیری با اطلاعات ممتاز (\en{OPSD}):} ژائو و همکاران \cite{zhao2026self} نشان دادند که بدون نیاز به آموزگار خارجی، می‌توان با شرطی‌سازی یک مدل بر روی پاسخ صحیح ($y^*$) یک «آموزگار درونی» ساخت که خط‌مشی بدون پاسخ دانش‌آموز را در مسیر استنتاج اصلاح کند:
    \begin{equation}
        \mathcal{L}_{\text{OPSD}}(\theta) = \mathbb{E}_{(x, y^*) \sim \mathcal{S}} \mathbb{E}_{\hat{y} \sim p_\theta(\cdot \mid x)} \left[ \sum_{t=1}^{|\hat{y}|} \mathcal{D}\left(p_\theta(\cdot \mid x, y^*, \hat{y}_{<t}) \,\|\, p_\theta(\cdot \mid x, \hat{y}_{<t})\right) \right]
    \end{equation}
    این رویکرد به ویژه در پایدارسازی آموزش مدل‌های زبانی در حل مسائل ریاضیات المپیادی نقشی کلیدی ایفا می‌کند \cite{song2026survey}.
\end{itemize}

\paragraph*{پایدارسازی پویایی آموزش و گزینش بهینه سختی (\en{PACED})}
یک نقیصه خطرناک در \en{OPD} بروز **تله پیشوند معیوب (\en{Flawed Prefix Trap})** است؛ بدین معنا که اگر دانش‌آموز توکنی غلط تولید کند، آموزگار روی بافت غلط شرطی شده و گرادیان‌های نویزی مخربی تولید می‌کند \cite{fu2026revisiting}. 
روش \en{PACED} \cite{xu2026paced} برای حل این مسئله، سختی پرومپت‌ها را بر اساس احتمال موفقیت دانش‌آموز ($p$) مدل‌سازی می‌کند. نسبت سیگنال به نویز (\en{SNR}) گرادیان تحت توزیع بتا بررسی می‌شود:
\begin{equation}
    w(p) = p^\alpha (1 - p)^\beta
\end{equation}
در حالت متقارن $\alpha = \beta = 1$، بیشینه گرادیان در نقطه مرزی توانمندی مدل ($p^* = 0.5$) محقق می‌شود و نمونه‌هایی که مدل در آن‌ها نرخ موفقیت صفر یا یک دارد (بسیار سخت یا بسیار آسان) به دلیل فقدان سیگنال گرادیان کارآمد فیلتر می‌گردند \cite{song2026survey}.

\paragraph*{همگرایی تقطیر دانش (\en{KD}) و تقطیر مجموعه داده (\en{Dataset Distillation - DD})}
همان‌طور که در پیمایش فانگ و همکاران \cite{fang2026knowledge} تبیین شده است، تقطیر داده‌ها به دنبال فشرده‌سازی مجموعه آموزشی عظیم $\mathcal{D}$ به یک مجموعه مصنوعی فشرده $\mathcal{D}_{\text{syn}}$ است. روش **\en{SRe2L}** \cite{yin2023squeeze} با تلفیق این دو مفهوم، از آموزگار برای بازبرچسب‌زنی احتمالات نرم داده‌های مصنوعی حاصل از وارون‌سازی مدل (\en{Model Inversion}) استفاده کرده و دانش‌آموز را بر روی این داده‌های بهینه‌شده آموزش می‌دهد که نیاز به ذخیره‌سازی داده‌های اصلی را تا ۹۰٪ کاهش می‌دهد.

\paragraph*{پایه‌های هندسی و نظری کوانتیزاسیون مدل‌های زبانی}
در زمینه تقطیر وزن‌ها برای استنتاج با حافظه اندک، چن و همکاران در مقاله \model{GPTQ-Babai} \cite{chen2026geometry} برای نخستین بار ساختار هندسی الگوریتم \en{GPTQ} \cite{frantar2023gptq} را رمزگشایی کردند.

\textbf{تبدیل به مسئله نزدیک‌ترین بردار (\en{Closest Vector Problem - CVP}):} مسئله کوانتیزاسیون لایه‌ای با تابع هدف خطای $L_2$ خروجی به صورت زیر است:
\begin{equation}
    \min_{z \in \mathbb{Z}^c} \|X \text{diag}(s) z - X w\|_2^2
\end{equation}
که در آن $X \in \mathbb{R}^{n \times c}$ ماتریس فعال‌سازی کالیبراسیون، $s \in \mathbb{R}_{>0}^c$ بردار مقیاس‌های کوانتیزاسیون، $w \in \mathbb{R}^c$ وزن اولیه و $z \in \mathbb{Z}^c$ بردار صحیح کوانتیزه‌شده است. این مسئله دقیقاً معادل با یافتن نزدیک‌ترین نقطه در یک شبکه لتیس با ماتریس پایه $B = X \text{diag}(s)$ به بردار هدف پیوسته $y = X w$ است:
\begin{equation}
    \min_{z \in \mathbb{Z}^c} \|B z - y\|_2^2
\end{equation}

\textbf{اثبات هم‌ارزی با الگوریتم نزدیک‌ترین صفحه بابای (\en{Babai's Algorithm}):} چن و همکاران \cite{chen2026geometry} اثبات کردند که گام به‌روزرسانی وزن‌ها در \en{OBQ} و \en{GPTQ} برای بعد $j_1$ پس از رند کردن بعد $j_2$:
\begin{equation}
    \Delta w[j_1] = - \frac{(B^\top B)^{-1}[j_1, j_2]}{(B^\top B)^{-1}[j_2, j_2]} (q[j_2] - w[j_2])
\end{equation}
دقیقاً معادل با تصویر ارتوگونال بردار باقی‌مانده هدف بر روی ابرصفحه تعریف‌شده توسط بردارهای متعامد گرام-اشمیت (\en{GSO}) در الگوریتم بابای \cite{babai1986lovasz} است. بر این مبنا، کران خطای لایه‌ای \en{GPTQ} در حالت بدون برش (\en{No-Clipping}) بر پایه تجزیه \en{LDL} ماتریس هسین $H = T^\top X^\top X T = L D L^\top$ به صورت تحلیلی و صلب استخراج شد:
\begin{equation}
    \|X \text{diag}(s) z - X w\|_2^2 \le \frac{1}{4} (T^{-1} s)^\top D (T^{-1} s)
\end{equation}

\paragraph*{انتشار خطای کوانتیزاسیون در شبکه‌های چندلایه (\en{QEP})}
آرای و ایچیکاوا در چارچوب \model{QEP} \cite{arai2025quantization} اثبات کردند که خطای کوانتیزاسیون در طول لایه‌های متوالی ترنسفورمر بدون کنترل به صورت **نمایی** منبسط می‌شود. آن‌ها با فرمول‌بندی تصحیح وزن تحلیلی:
\begin{equation}
    \widetilde{W}_l^* = W_l + W_l \delta_l \widehat{X}_l^\top \widehat{H}_l^{-1}
\end{equation}
که در آن $\delta_l = X_l - \widehat{X}_l$ عدم انطباق فعال‌سازی‌های دقت کامل و کوانتیزه‌شده است، نشان دادند که این بهینه‌سازی دقیقاً معادل با رگولاریزاسیون ریج (\en{Ridge Regularization}) برای مهار بیش‌برازش روی داده‌های کالیبراسیون است:
\begin{equation}
    \min_{\widetilde{W}_l} \|W_l X_l - \widetilde{W}_l \widehat{X}_l\|_F^2 + \lambda_l \|W_l - \widetilde{W}_l\|_F^2
\end{equation}

\paragraph*{ناهمدوسی هادامارد و سرکوب داده‌های پرت}
پژوهش‌های اخیر نشان داده‌اند که فعال‌سازی‌های مدل‌های زبانی به شدت دارای داده‌های پرت شدید (\en{Outliers}) هستند که کوانتیزاسیون به زیر ۴ بیت را ناممکن می‌سازد \cite{dettmers2022gpt3, maisonnave2025gradual, panferov2025quest}. میسوناو و همکاران \cite{maisonnave2025gradual} به صورت نظری نشان دادند که اعمال ماتریس تبدیل هادامارد (\en{Hadamard Matrix}) به مراتب بهتر از ماتریس‌های ارتوگونال تصادفی داده‌های پرت را سرکوب می‌کند. 
به ازای بردار $x = (c, \epsilon, \dots, \epsilon)^\top$ با مقدار فرین $c \gg \epsilon$:
\begin{equation}
    \max_{i} |(H x)_i| = \frac{c}{\sqrt{n}} \quad \text{vs.} \quad \mathbb{E}\left[\max_{i} |(Q x)_i|\right] = c \sqrt{\frac{2 \log n}{n}}
\end{equation}
علاوه بر این، آن‌ها اثبات کردند که ماتریس‌های هادامارد با رسیدن به کران پایین $\max_{i,j} |Q_{ij}| \ge \frac{1}{\sqrt{n}}$ بهینه‌ترین کانال بازتوزیع انرژی بردار در فضای چندبعدی به شمار می‌روند \cite{maisonnave2025gradual}.

\paragraph*{قضیه خطی‌بودن پرپلکسیتی (\en{The Linearity Theorem})}
مالینوفسکی و همکاران \cite{malinovskii2024pushing} ارتباط مستقیم بین خطای میانگین مربعات ($L_2$ \en{MSE}) لایه‌ها و معیار پرپلکسیتی کل مدل را با بسط تیلور مرتبه دوم اثبات نمودند:
\begin{equation}
    \mathbb{E}[\text{PPL}(\widehat{W})] \approx \text{PPL}(W^*) + \sum_{l=1}^L \alpha_l t_l^2, \quad t_l^2 = \frac{\mathbb{E}\|W_l - \widehat{W}_l\|_F^2}{\|W_l\|_F^2}
\end{equation}
که در آن $\alpha_l$ ضرایب حساسیت مستقل از نوع فشرده‌سازی هستند. این فرمول‌بندی پایه الگوریتم **\en{HIGGS}** گردید که با ترکیب چرخش‌های هادامارد و گریدهای بهینه گاوسی، رکوردهای جدیدی در فشرده‌سازی بدون داده ثبت کرد \cite{malinovskii2024pushing}.

\paragraph*{کوانتیزاسیون بدون اتلاف آماری (\en{SLQ}) و قانون واریانس گاما-توان‌دو}
در پژوهش هلسیگ و همکاران \cite{helcig2026statistically}، مرز میان فشرده‌سازی اتلافی و بدون اتلاف با معرفی مفهوم **کوانتیزاسیون بدون اتلاف آماری (\en{Statistically-Lossless Quantization - SLQ})** بازتعریف شد. آن‌ها دو معیار ارائه دادند:
\begin{itemize}
    \item \textbf{اتلاف‌ناپذیری در سطح وظیفه (\en{Task-Lossless - TL}):} حفظ دقت بنچمارک‌ها در دامنه واریانس نمونه‌برداری استوکستیک ($\ge 99\%$).
    \item \textbf{اتلاف‌ناپذیری در سطح توزیع (\en{Distribution-Lossless - DL}):} انطباق کامل توزیع توکن بعدی که با معیار تحلیلی **نرخ پذیرش مورد انتظار (\en{Expected Acceptance Rate - EAR})** سنجیده می‌شود:
    \begin{equation}
        \text{EAR} = 1 - d_{\text{TV}}(p, q) = \sum_{k=1}^{|\mathcal{V}|} \min(p(k), q(k))
    \end{equation}
\end{itemize}

\textbf{استخراج قانون واریانس گاما-توان‌دو (\en{Gamma-Squared Variance Law}):}
برای بردار وزن‌ها در بازه $[L, U]$ با دامنه $R = U - L$ و حداکثر اندازه $M = \max(|L|, |U|)$، شاخص ناکارآمدی مرکزی (\en{Centering Inefficiency}) به فرم $\gamma = \frac{2M}{R} \ge 1$ تعریف می‌شود. اندازه‌های گام کوانتیزاسیون متقارن و نامتقارن عبارتند از:
\begin{equation}
    \Delta_{\text{asym}} = \frac{R}{n-1}, \quad \Delta_{\text{sym}} = \frac{2M}{n-1} \implies \Delta_{\text{sym}} = \gamma \Delta_{\text{asym}}
\end{equation}
تحت تقریب نرخ‌بالای بنت \cite{bennett1948spectra} که واریانس نویز یکنواخت برابر با $\sigma^2 = \frac{\Delta^2}{12}$ است، رابطه مستقیم زیر استخراج می‌شود \cite{helcig2026statistically}:
\begin{equation}
    \sigma_{\text{sym}}^2 = \frac{\Delta_{\text{sym}}^2}{12} = \frac{(\gamma \Delta_{\text{asym}})^2}{12} = \gamma^2 \sigma_{\text{asym}}^2
\end{equation}
این نتیجه اثبات می‌کند که در مدل‌های زبانی نامتقارن، کوانتیزاسیون متقارن واریانس نویز را با ضریب $\gamma^2$ تشدید کرده و به شکستن تصمیمات حاشیه‌ای سافت‌مکس می‌انجامد؛ بنابراین، **کوانتیزاسیون نامتقارن پیش‌شرط قطعی برای فشرده‌سازی بدون اتلاف توزیعی است**.

\paragraph*{رویکردهای مکمل: \model{ResQ} و \model{QuEST}}
\begin{itemize}
    \item \textbf{روش \model{ResQ} \cite{saxena2025resq}:} بهره‌گیری از تحلیل مؤلفه‌های اصلی (\en{PCA}) بر روی ماتریس کوواریانس فعال‌سازی‌ها ($XX^\top$) جهت شناسایی زیرفضای کم‌رتبه ($r = d/8$) با واریانس بالا و نگهداری آن در دقت ۸ بیتی، در حالی که مابقی درایه‌ها با چرخش تصادفی در ۴ بیت کوانتیزه می‌شوند. کران خطای اثبات‌شده تضمین می‌کند که بیشینه‌سازی انرژی در مؤلفه‌های اصلی، خطای فروبنیوس را به حداقل می‌رساند.
    \item \textbf{روش \model{QuEST} \cite{panferov2025quest}:} ارائه نخستین راهکار آموزش آگاه از کوانتیزاسیون (\en{QAT}) پایدار برای مدل‌های ۱ بیتی و ۲ بیتی با تعریف برآوردگر گرادیان مطمئن (\en{Trust Gradient Estimator}) بر پایه ماسک‌گذاری نقاط با خطای کوانتیزاسیون بالا در حوزه تبدیل هادامارد:
    \begin{equation}
        \frac{\partial}{\partial x} \approx \text{IHT}\left( M_{\alpha^*}(x_h; \hat{x}_h) \odot \frac{\partial}{\partial \hat{x}_h} \right), \quad M_{\alpha^*} = \mathbb{I}_{|\hat{x} - x| \le T}
    \end{equation}
\end{itemize}

\paragraph*{ارزیابی‌های آماری و تجربی روش‌های فشرده‌سازی و تقطیر}
در جدول زیر، مقایسه جامع عملکرد، سطوح فشرده‌سازی و میزان حفظ قابلیت‌های استدلالی در روش‌های شاخص مرور شده بر پایه نتایج رسمی مقالات مرجع گزارش شده است:

\begin{table}[htbp]
\centering
\caption{مقایسه عملکردی روش‌های نوین تقطیر و کوانتیزاسیون در مدل‌های زبانی بزرگ}
\label{tab:kd_quant_sota_eval}
\resizebox{\textwidth}{!}{%
\begin{tabular}{llccccc}
\toprule
\textbf{روش} & \textbf{پارادایم / رویکرد} & \textbf{پیکربندی بیت (\en{W/A/KV})} & \textbf{مدل ارزیابی} & \textbf{پرپلکسیتی (\en{WikiText-2})} & \textbf{دقت بنچمارک هدف} & \textbf{شتاب استنتاج} \\
\midrule
\textbf{\en{BF16 Baseline}} & دقت پایه & 16 / 16 / 16 & \model{Qwen3-8B} & 9.73 & 73.02\% (\en{MMLU}) & 1.0× \\
\textbf{\en{FP16 Baseline}} & دقت پایه & 16 / 16 / 16 & \model{Llama-3.1-8B} & 7.20 & 68.31\% (\en{MMLU}) & 1.0× \\
\midrule
\model{GPTQ} \cite{frantar2023gptq} & \en{PTQ} (حریصانه / \en{Hessian}) & 4 / 16 / 16 & \model{Qwen3-8B} & 10.10 & 71.76\% (\en{MMLU}) & 2.1× \\
\model{QEP-GPTQ} \cite{arai2025quantization} & \en{PTQ} (تصحیح خطای انباشته) & 4 / 16 / 16 & \model{Llama-2-7B} & 5.93 & 67.95\% (\en{Avg-3Task}) & 2.1× \\
\model{QEP-QuIP} \cite{arai2025quantization} & \en{PTQ} (چرخش هادامارد + \en{QEP}) & 2 / 16 / 16 & \model{Llama-2-70B} & 5.86 & 69.98\% (\en{Avg-3Task}) & 3.8× \\
\model{HIGGS} \cite{malinovskii2024pushing} & \en{PTQ} (قضیه خطی‌بودن + گرید گاوسی) & 4.02 / 16 / 16 & \model{Llama-3.1-8B} & 6.01 & 63.26\% (\en{MMLU}) & 2.5× \\
\model{HPTQ} \cite{chen2026geometry} & لتیس / بابای بدون برش & 3.125 / 16 / 16 & \model{Qwen3-8B} & 10.34 & 70.96\% (\en{MMLU}) & 2.4× \\
\model{SSQR-1\%} \cite{chen2026geometry} & تطبیق مقیاس + پسماند اسپارس & 3.445 / 16 / 16 & \model{Qwen3-8B} & 10.64 & 68.46\% (\en{MMLU}) & 2.0× \\
\model{ResQ} \cite{saxena2025resq} & تجزیه \en{PCA} کم‌رتبه + چرخش & 4 / 4 / 4 (1/8 INT8) & \model{Llama-3-8B} & 7.10 & 57.20\% (\en{MMLU}) & 3.03× \\
\model{QuEST} \cite{panferov2025quest} & \en{QAT} (برآوردگر گرادیان مطمئن) & 4 / 4 / 16 & \model{Llama-800M} & 2.65 (\en{C4 Loss}) & 71.16\% (\en{PiQA}) & 2.2× \\
\model{SLQ-TL} \cite{helcig2026statistically} & تخصیص پویا شپلی (اتلاف‌ناپذیر وظیفه) & 3.59 / 16 / 16 & \model{Llama-3.3-70B} & 9.80 & 99.1\% (\en{Recovery}) & 3.68× \\
\model{SLQ-DL} \cite{helcig2026statistically} & تخصیص پویا شپلی (اتلاف‌ناپذیر توزیع) & 5.70 / 16 / 16 & \model{Qwen3-8B} & 9.81 (\en{EAR: 99\%}) & 99.3\% (\en{Recovery}) & 1.74× \\
\midrule
\model{OPSD} \cite{zhao2026self} & خود-تقطیری درون‌خط‌مشی با داده ممتاز & 16 / 16 / 16 & \model{Qwen3-8B-Inst} & - & 71.4\% (\en{AIME/HMMT}) & 8× کاهش سمپل \\
\bottomrule
\end{tabular}%
}
\end{table}

\paragraph*{چالش‌های حل‌نشده، پاتولوژی‌های شکست و چشم‌انداز آینده}
علیرغم پیشرفت‌های حاصله، چندین چالش بنیادین در کانون توجه پژوهشگران قرار دارد \cite{song2026survey, fang2026knowledge}:
\begin{itemize}
    \item \textbf{اشباع در خود-تقطیری و مسئله اوروبوروس (\en{The Ouroboros Problem}):} در خود-تقطیری بدون بازخورد محیطی یا برچسب مرجع، مدل خطاهای ساختاری خود را بازتولید و تقویت کرده که منجر به فروپاشی تنوع توزیع می‌گردد \cite{song2026survey}.
    \item \textbf{شکاف قابلیت-کالیبراسیون (\en{Calibration-Capability Gap}):} مدل‌های تقطیرشده با \en{OPD} علیرغم کسب نمرات بالا در حل مسائل، دچار بیش‌اعتمادی شدید شده و توانایی خودارزیابی عدم قطعیت (\en{Epistemic Uncertainty}) را از دست می‌دهند \cite{zhang2026illusion}.
    \item \textbf{تورم توکن‌های تفکر (\en{Thinking Token Inflation}):} فشرده‌سازی‌های اتلافی ناخواسته منجر به طولانی‌تر شدن ردپای استدلال مدل (تا ۲.۵ برابر) برای رسیدن به پاسخ یکسان می‌شوند که مزیت کارایی حافظه را تضعیف می‌سازد \cite{helcig2026statistically}.
    \item \textbf{تدوین قوانین مقیاس‌پذیری تقطیر (\en{Distillation Scaling Laws}):} استخراج روابط توانی دقیق مشابه قوانین کاپلان و چینچیلا برای تخصیص بهینه بودجه میان اندازه آموزگار، اندازه دانش‌آموز و میزان رول‌آوت‌های درون‌خط‌مشی از مهم‌ترین مسائل باز در این حوزه است \cite{busbridge2025distillation, song2026survey}.
\end{itemize}